% !TeX root = ../main.tex
\chapter{Insiemi finiti, infiniti e cardinalità}
\fonti{Lezioni 7, 8 e 9}

\section{Insiemi finiti e infiniti}
% A–L L7 p. 2; M–Z L8 pp. 1–2. Le proprietà elementari finite sono assunte.
\begin{definizione}
Poniamo $[n]=\{1,\ldots,n\}$ per $n\ge1$ e $[0]=\varnothing$. Un insieme $X$ è \textbf{finito} se esiste una biezione $[n]\to X$ per qualche $n\in\N$; è \textbf{infinito} se non è finito. Nel caso finito il numero $n$, univocamente determinato, si chiama \textbf{cardinalità} e si indica con $|X|$.
\end{definizione}

Assumiamo i fatti elementari sul conteggio finito: $[m]$ e $[n]$ sono in biezione solo se $m=n$; un sottoinsieme di un insieme finito è finito; un'iniezione di un insieme finito in se stesso è suriettiva. In particolare un insieme finito non può essere in biezione con un proprio sottoinsieme stretto.

\begin{corollario}
Se $X$ è finito, $|\mathcal P(X)|=2^{|X|}$.
\end{corollario}
\begin{proof}
Per il teorema~\ref{teo:caratteristiche}, i sottoinsiemi di $X$ corrispondono alle funzioni $X\to\{0,1\}$. Se $|X|=n$, per ciascuno degli $n$ elementi si sceglie indipendentemente uno dei due valori: le scelte sono $2^n$. Per $n=0$ c'è la sola funzione vuota e $\mathcal P(\varnothing)=\{\varnothing\}$ ha $1=2^0$ elemento.
\end{proof}

\begin{teorema}\label{teo:infiniti}
Nel quadro con assioma della scelta, per un insieme $X$ sono equivalenti:
\begin{enumerate}[label=\arabic*.]
\item $X$ è infinito;
\item esiste un'iniezione $f:\N\to X$;
\item esiste un'iniezione $g:X\to X$ non suriettiva;
\item esiste una biezione fra $X$ e un suo sottoinsieme proprio.
\end{enumerate}
\end{teorema}
\begin{proof}
% A–L L7 pp. 2–3 e L8 pp. 1–2, più dettagliata sui tre casi dello spostamento.
\passomath{1\Rightarrow2} Poiché $X$ è infinito, è non vuoto: scegliamo $x_0\in X$. Scelti $x_0,\ldots,x_{n-1}$ distinti, il complementare $X\setminus\{x_0,\ldots,x_{n-1}\}$ è non vuoto, altrimenti $X$ sarebbe finito. Possiamo dunque scegliere $x_n$ in esso. L'assioma della scelta garantisce la successione di scelte per tutti gli $n$. La funzione $f(n)=x_n$ è iniettiva: per $m<n$, $x_n$ è stato scelto diverso da $x_m$.

\passomath{2\Rightarrow3} Posto $S=f(\N)$, definiamo
\[
g(x)=\begin{cases}f(n+1)&x=f(n)\in S,\\x&x\notin S.\end{cases}
\]
L'indice $n$ è unico perché $f$ è iniettiva; la regola è quindi ben definita. Per l'iniettività distinguiamo i casi. Se $x=f(n)$ e $x'=f(m)$, da $g(x)=g(x')$ segue $f(n+1)=f(m+1)$, dunque $n=m$ e $x=x'$. Se entrambi sono fuori da $S$, sono lasciati fissi, quindi immagini uguali impongono punti uguali. Se solo uno è in $S$, anche solo una delle immagini è in $S$: non possono coincidere. Infine $f(0)$ non è raggiunto: gli elementi fuori da $S$ rimangono fuori, e quelli in $S$ vengono mandati nei punti $f(n+1)\ne f(0)$. Dunque $g$ non è suriettiva.

\passomath{3\Rightarrow4} Restringendo il codominio di $g$ a $g(X)$ si ottiene una biezione $X\to g(X)$. La mancata suriettività dice esattamente $g(X)\subsetneq X$.

\passomath{4\Rightarrow1} Una biezione $h:X\to S\subsetneq X$, composta con l'inclusione $S\hookrightarrow X$, dà un'iniezione non suriettiva $X\to X$. Questo è impossibile per un insieme finito, per il fatto di conteggio assunto sopra. Quindi $X$ è infinito.
\end{proof}

\begin{schema}
\begin{tikzpicture}[>=Stealth]
  \foreach \i in {0,...,3}{\node (x\i) at (2*\i,1) {$f(\i)$};\node (y\i) at (2*\i,-.5) {$f(\i)$};}
  \foreach \i in {0,...,2}{\pgfmathtruncatemacro{\j}{\i+1}\draw[->,coloreschema] (x\i)--(y\j);}
  \node at (8,1) {$\cdots$};\node at (8,-.5) {$\cdots$};\draw[->,coloreschema] (x3)--(7.7,-.3);
  \node[below=2pt] at (y0.south) {omesso};\node at (4,-1.6) {fuori da $f(\N)$ ogni elemento resta fisso};
\end{tikzpicture}
\end{schema}

\begin{esempio}
L'applicazione $n\mapsto n+1$ è un'iniezione $\N\to\N$ che non raggiunge $0$; cambiando il codominio dà una biezione $\N\to\N^*$. Anche $n\mapsto2n$ dà una biezione fra $\N$ e il sottoinsieme proprio $2\N$ dei naturali pari. Sono esempi concreti di una possibilità esclusa per gli insiemi finiti.
\end{esempio}

\section{Equipotenza e confronto di cardinalità}
% A–L L8 pp. 2–3; M–Z L8 pp. 5–6.
\begin{definizione}
Due insiemi sono \textbf{equipotenti}, e si scrive $X\simeq Y$ oppure $|X|=|Y|$, se esiste una biezione $X\to Y$. Si scrive $|X|\le|Y|$ se esiste un'iniezione $X\to Y$; si scrive $|X|<|Y|$ se vale $\le$ ma non l'uguaglianza.
\end{definizione}

L'equipotenza è riflessiva, tramite l'identità; simmetrica, tramite l'inversa di una biezione; transitiva, tramite la composizione di biezioni. La notazione $|X|$ per insiemi arbitrari esprime la grandezza comune agli insiemi equipotenti, senza richiedere qui una costruzione formale dei cardinali. Per esempio $|\{\text{rosso},\text{blu},\text{verde}\}|=3$.

Per insiemi non vuoti, $|X|\le|Y|$ equivale all'esistenza di una suriezione $Y\to X$, per il corollario sui criteri delle inverse. La formulazione tramite iniezioni vale anche per il vuoto: $|\varnothing|\le|Y|$ sempre, benché non esista una funzione $Y\to\varnothing$ se $Y\ne\varnothing$.

\begin{teorema}[Cantor--Schröder--Bernstein]\label{teo:csb}
Se esistono iniezioni $X\to Y$ e $Y\to X$, allora esiste una biezione $X\to Y$. Dunque $|X|\le|Y|$ e $|Y|\le|X|$ implicano $|X|=|Y|$.
\end{teorema}
Il teorema è assunto senza dimostrazione. Il confronto delle cardinalità è riflessivo e transitivo perché l'identità è iniettiva e la composizione di iniezioni è iniettiva; l'antisimmetria è precisamente il risultato appena enunciato.

\begin{teorema}[Confronto totale]
Nel quadro con assioma della scelta, per ogni coppia di insiemi $X,Y$ vale $|X|\le|Y|$ oppure $|Y|\le|X|$.
\end{teorema}
Anche questo risultato è assunto senza dimostrazione. In particolare, per il teorema~\ref{teo:infiniti}, la cardinalità di $\N$ è la più piccola cardinalità infinita.

\section{Numerabilità}
\begin{definizione}
Un insieme $X$ è \textbf{numerabile} se $|X|\le|\N|$; è \textbf{numerabile infinito} se $|X|=|\N|$. La convenzione include fra i numerabili gli insiemi finiti e il vuoto. Scriviamo $\aleph_0=|\N|$.
\end{definizione}

Un insieme numerabile infinito è in biezione con $\N$: se $X\hookrightarrow\N$ e $X$ è infinito, il teorema~\ref{teo:infiniti} dà anche $\N\hookrightarrow X$, e si applica Cantor--Schröder--Bernstein.

\begin{proposizione}
\begin{enumerate}[label=\arabic*.]
\item $|\N\times\N|=\aleph_0$.
\item $|X^r|=\aleph_0$ per ogni insieme $X$ con $|X|=\aleph_0$ e ogni intero $r\ge1$; in particolare $|\N^r|=\aleph_0$.
\end{enumerate}
\end{proposizione}
\begin{proof}
\begin{enumerate}[label=\arabic*.]
\item \textbf{Coppie.} Consideriamo $p:\N\times\N\to\N$, $p(m,n)=2^m(2n+1)-1$. Se $p(m,n)=p(m',n')$, allora
\[
2^m(2n+1)=2^{m'}(2n'+1).
\]
Se $m<m'$, cancellando $2^m$ si avrebbe un numero dispari uguale a uno pari, impossibile. Lo stesso esclude $m>m'$ scambiando i ruoli. Quindi $m=m'$, e cancellando la potenza di $2$ si ricava $n=n'$. L'applicazione è iniettiva; nell'altro verso $m\mapsto(m,0)$ è un'iniezione. Cantor--Schröder--Bernstein dà l'equipotenza. Non si è usata la fattorizzazione unica in primi.
\item \textbf{Prodotti finiti, induzione su $r$.} Per $r=1$ vale l'identità. Supponiamo data una biezione $u_r:\N^r\to\N$. L'applicazione $(a_1,\ldots,a_{r+1})\mapsto(u_r(a_1,\ldots,a_r),a_{r+1})$ è una biezione $\N^{r+1}\to\N^2$: l'inversa applica $u_r^{-1}$ alla prima coordinata. Componendola con una biezione $\N^2\to\N$ si ottiene il passo $r+1$. Il principio d'induzione conclude. Per un insieme $X$ con $|X|=\aleph_0$, una biezione $v:\N\to X$ induce la biezione $(n_1,\ldots,n_r)\mapsto(v(n_1),\ldots,v(n_r))$ da $\N^r$ a $X^r$, con inversa coordinata per coordinata. Quindi anche $|X^r|=\aleph_0$.
\end{enumerate}
\end{proof}

\begin{osservazione}[Enumerazione per diagonali]
Un secondo metodo consiste nel percorrere successivamente le diagonali $m+n=0,1,2,\ldots$. Ciascuna ha un numero finito di punti e ogni coppia appartiene a una sola diagonale. Elencando ogni diagonale una volta si ottiene una successione senza ripetizioni che raggiunge tutte le coppie. Questo dà direttamente una biezione $\N\to\N^2$.
\end{osservazione}

\begin{schema}
\begin{tikzpicture}[scale=.55,>=Stealth]
  \draw[asse] (-.4,0)--(5.5,0) node[right] {$m$};\draw[asse] (0,-.4)--(0,5.5) node[above] {$n$};
  \foreach \x in {0,...,4}{\foreach \y in {0,...,4}{\fill (\x,\y) circle (1.7pt);}}
  \draw[->,coloreschema,thick] (0,0)--(1,0)--(0,1)--(0,2)--(1,1)--(2,0)--(3,0)--(2,1)--(1,2)--(0,3);
  \node[right,align=left] at (6,2.5) {Un'altra enumerazione percorre le diagonali\\$m+n=0,1,2,\ldots$ del reticolo.};
\end{tikzpicture}
\end{schema}

\begin{teorema}[Unione numerabile di numerabili]\label{teo:unione-numerabile}
Nel quadro con assioma della scelta, siano $I$ numerabile e $A_i$ numerabile per ogni $i\in I$.
\begin{enumerate}[label=\arabic*.]
\item $\bigcup_{i\in I}A_i$ è numerabile.
\item Se almeno un $A_j$ è infinito, l'unione è numerabile infinita: $|\bigcup_{i\in I}A_i|=\aleph_0$.
\end{enumerate}
\end{teorema}
\begin{proof}
\textbf{1. Riduzione ai membri non vuoti.} Se l'unione è vuota, la conclusione vale. Altrimenti poniamo $I'=\{i\in I\mid A_i\ne\varnothing\}$. È numerabile e non vuoto, quindi esiste una suriezione $e:\N\to I'$: se $I'$ è infinito si usa una biezione, se è finito si ripete un elenco dei suoi elementi.

\textbf{Costruzione della suriezione.} Ogni $A_i$ non vuoto è a sua volta immagine di una suriezione $\varphi_i:\N\to A_i$. Scegliamo una tale funzione per ciascun $i\in I'$: è qui che occorre la scelta simultanea. Definiamo
\[
\Phi:\N^2\to\bigcup_{i\in I}A_i,\qquad \Phi(k,n)=\varphi_{e(k)}(n).
\]
Per ogni $a$ dell'unione, esiste $i\in I'$ con $a\in A_i$; esiste $k$ con $e(k)=i$ e poi $n$ con $\varphi_i(n)=a$. Quindi $\Phi$ è suriettiva. Componendo con una biezione $\N\to\N^2$ si ha una suriezione da $\N$ sull'unione, dunque un'iniezione dell'unione in $\N$ per il criterio già provato. L'unione è numerabile.

\textbf{2.} Se $A_j$ è infinito, l'inclusione $A_j\hookrightarrow\bigcup_i A_i$ e il teorema~\ref{teo:infiniti} danno un'iniezione di $\N$ nell'unione. Il punto~1 dà l'iniezione opposta verso $\N$, quindi l'uguaglianza delle cardinalità segue da Cantor--Schröder--Bernstein. In particolare si recupera il caso in cui $I$ e tutti gli $A_i$ sono numerabili infiniti.
\end{proof}

\begin{corollario}
\begin{enumerate}[label=\arabic*.]
\item $|\Z|=\aleph_0$.
\item $|\Q|=\aleph_0$.
\end{enumerate}
\end{corollario}
\begin{proof}
\begin{enumerate}[label=\arabic*.]
\item \textbf{Interi.} Le formule $e(2k)=k$ ed $e(2k+1)=-(k+1)$ definiscono una biezione $\N\to\Z$: i due rami raggiungono rispettivamente tutti i non negativi e tutti i negativi, senza ripetizioni.
\item \textbf{Razionali.} Il prodotto $\Z\times\N^*$ è numerabile perché si inietta in $\N^2$, usando una biezione $\Z\to\N$ e l'inclusione $\N^*\subseteq\N$. La funzione $(a,b)\mapsto a/b$ è suriettiva su $\Q$: ogni razionale ammette un denominatore positivo cambiando entrambi i segni, se necessario. Dunque $|\Q|\le\aleph_0$. L'iniezione $\N\hookrightarrow\Q$, $n\mapsto n/1$, dà la disuguaglianza opposta e si conclude per Cantor--Schröder--Bernstein.
\end{enumerate}
\end{proof}

\section{Cantor e la cardinalità del continuo}
% A–L L8 pp. 5–6; M–Z L8 pp. 2–3. Nessun anticipo di Cantor prima degli infiniti.
\begin{teorema}[Cantor]\label{teo:cantor}
Per ogni insieme $X$ non esiste una suriezione $X\to\mathcal P(X)$. In particolare $|X|<|\mathcal P(X)|$.
\end{teorema}
\begin{proof}
\textbf{Diagonale.} Data una qualunque $f:X\to\mathcal P(X)$, definiamo $D=\{x\in X\mid x\notin f(x)\}$. È un sottoinsieme di $X$, dunque un elemento del codominio. Se $D=f(a)$, allora
\[
a\in D\iff a\notin f(a)\iff a\notin D.
\]
Se $a\in D$ si ottiene $a\notin D$; se $a\notin D$ si ottiene $a\in D$. Entrambi i casi sono impossibili. Quindi $D$ non è nell'immagine e $f$ non è suriettiva.

\textbf{Confronto stretto.} La funzione $x\mapsto\{x\}$ è iniettiva, quindi $|X|\le|\mathcal P(X)|$. Non può esserci una biezione, perché sarebbe una suriezione. Per $X=\varnothing$, la conclusione vale direttamente: $\mathcal P(X)=\{\varnothing\}$ ha un elemento.
\end{proof}

\begin{definizione}
Si pone $2^{|X|}=|\{0,1\}^X|=|\mathcal P(X)|$. La \textbf{cardinalità del continuo} è $\mathfrak c=2^{\aleph_0}$. Il teorema di Cantor dà $\aleph_0<\mathfrak c$ e, iterato, dà
\[
\aleph_0<2^{\aleph_0}<2^{\,2^{\aleph_0}}<\cdots.
\]
\end{definizione}

\begin{osservazione}[Ipotesi del continuo]
L'\textbf{ipotesi del continuo} afferma che non esiste una cardinalità strettamente compresa fra $\aleph_0$ e $2^{\aleph_0}$. Non è una conseguenza del solo teorema di Cantor. Si assume qui, senza dimostrarlo, il risultato di logica secondo cui essa è indipendente dagli assiomi usuali degli insiemi con scelta, se questi sono coerenti. L'ipotesi stessa non viene assunta per i risultati di questa dispensa.
\end{osservazione}

\begin{esercizio}
Siano $A$ finito e $B$ numerabile. Sono numerabili $A\cup B$, $A\times B$ e $B^A$?
\end{esercizio}
\begin{svolgimento}
\textbf{Unione.} Un insieme finito è numerabile. Il teorema~\ref{teo:unione-numerabile}, applicato alla famiglia di due membri $A,B$, dà la numerabilità dell'unione.

\textbf{Prodotto.} Iniezioni $A\to\N$ e $B\to\N$ inducono, coordinata per coordinata, un'iniezione $A\times B\to\N^2$. Quindi il prodotto è numerabile.

\textbf{Funzioni.} Per $A=\{a_1,\ldots,a_m\}$, $m\ge1$, la mappa $f\mapsto(f(a_1),\ldots,f(a_m))$ è una biezione $B^A\to B^m$, perché ogni scelta della $m$-upla assegna esattamente una funzione. L'iniezione $B\to\N$ dà $B^m\hookrightarrow\N^m$, che è numerabile. Se $A=\varnothing$, $B^A$ ha la sola funzione vuota; se $B=\varnothing$ e $A\ne\varnothing$, non esistono funzioni e $B^A=\varnothing$. Anche nei casi vuoti tutte e tre le risposte sono affermative.
\end{svolgimento}

\section{Cardinalità della retta reale}
% A–L L9 p. 1; M–Z L9 pp. 2–3. Passo analitico assunto, come dichiarato in M–Z.
\begin{proposizione}
Una biezione $u:X\to Y$ induce una biezione $u_*:\mathcal P(X)\to\mathcal P(Y)$.
\end{proposizione}
\begin{proof}
La candidata inversa è $(u^{-1})_*$. Le regole di composizione delle immagini dirette danno $(u^{-1})_*\circ u_*=(u^{-1}\circ u)_*=(\operatorname{id}_X)_*$, che è l'identità di $\mathcal P(X)$. Nell'altro verso, $u_*\circ(u^{-1})_*=(\operatorname{id}_Y)_*$, identità di $\mathcal P(Y)$. Dunque le due applicazioni sono inverse.
\end{proof}

\begin{teorema}\label{teo:continuo}
$|\R|=\mathfrak c=2^{\aleph_0}$.
\end{teorema}
\begin{proof}[Dimostrazione con un risultato analitico assunto]
\passomath{|\R|\le\mathfrak c} A ogni $x\in\R$ associamo il taglio razionale inferiore $L_x=\{q\in\Q\mid q<x\}$. Se $x<y$, per densità esiste $q\in\Q$ con $x<q<y$: esso appartiene a $L_y$, ma non a $L_x$. La mappa $x\mapsto L_x$ è quindi iniettiva. Poiché $\Q\simeq\N$, la proposizione precedente dà $|\mathcal P(\Q)|=|\mathcal P(\N)|=\mathfrak c$, dunque $|\R|\le\mathfrak c$.

\passomath{\mathfrak c\le|\R|} Si assume il risultato analitico che esiste un'iniezione $\{0,1\}^{\N}\to\R$.\footnote{La costruzione mediante serie numeriche richiede un argomento aggiuntivo non sviluppato qui. La formula binaria da sola non prova l'iniettività: $(1,0,0,\ldots)$ e $(0,1,1,\ldots)$ danno entrambe $1$ nella somma $\sum_{n\ge0}\varepsilon(n)/2^n$.} Da questo fatto e dalla prima disuguaglianza segue l'uguaglianza per Cantor--Schröder--Bernstein.
\end{proof}

\section{Somme e prodotti di cardinalità infinite}
% Integrazione M–Z L9 pp. 4–5, espressamente autorizzata dall'autore.
\begin{definizione}
L'\textbf{unione disgiunta} è $A\sqcup B=(A\times\{0\})\cup(B\times\{1\})$: le etichette mantengono distinte le copie anche se $A\cap B\ne\varnothing$. Per $|A|=\alpha$, $|B|=\beta$, si definiscono $\alpha+\beta=|A\sqcup B|$ e $\alpha\beta=|A\times B|$.
\end{definizione}

Le definizioni non dipendono dagli insiemi che rappresentano le cardinalità. Infatti biezioni $u:A\to A'$ e $v:B\to B'$ inducono, sul prodotto, $(a,b)\mapsto(u(a),v(b))$, con inversa coordinata per coordinata; sull'unione disgiunta si usano $(a,0)\mapsto(u(a),0)$ e $(b,1)\mapsto(v(b),1)$, ancora invertibili.

\begin{teorema}
Nel quadro con assioma della scelta, se $\alpha,\beta$ sono cardinalità non nulle e almeno una è infinita,
\[
\alpha+\beta=\alpha\beta=\max\{\alpha,\beta\}.
\]
\end{teorema}
Il risultato generale è assunto senza dimostrazione. L'ipotesi di fattori non vuoti è necessaria per il prodotto: se $B=\varnothing$, $A\sqcup B\simeq A$, ma $A\times B=\varnothing$.
